% Definiciones y constantes de estilo
\include{definiciones}

% Definiciones de comandos
\newcommand{\nombreautor}{Javier López Cano}
\newcommand{\nombretutor}{Ramón Díaz Uriarte}
\newcommand{\nombretrabajo}{TODO: Título del Trabajo de Fin de Grado}
\newcommand{\fecha}{Junio 2021}
\newcommand{\grado}{Grado en Ingeniería Informática}
% Descomentar si tu trabajo tiene un ponente
%\newcommand{\nombreponente}{TODO: Nombre del ponente}
% Descomentar si tu trabajo está asociado a un grupo de investigación
% \newcommand{\grupoInvestigacion}{TODO: Grupo de investigación}
\newcommand{\departamento}{Dpto. de Bioquímica}
\newcommand{\facultad}{Escuela Politécnica Superior}
\newcommand{\universidad}{Universidad Autónoma de Madrid}
\newcommand{\pieparizq}{TODO: Pie de página par}
\newcommand{\pieparcen}{Trabajo de Fin de Grado}
\newcommand{\logoizq}{Logo_EPS}
\newcommand{\logoder}{Logo_UAM}
\newcommand{\correo}{javier.lopezcano@estudiante.uam.es}

% Glosario y acrónimos
\makeglossaries

% Glosario

\newglossaryentry{cancer}{name={cáncer},description={Conjunto de enfermedades caracterizadas por que un conjunto de células comienza a reproducirse de forma muy rápida produciendo un crecimiento anormal, en ocasiones además pueden desplazarse a otros tejidos del cuerpo},plural={cánceres}}
\newglossaryentry{tumor}{name={tumor},description={Conjunto de células anormalmente grande que se produce por la multiplicación más rápida de lo normal de un conjunto de células. En ocasiones son producidos por cáncer, en cuyo caso pueden invadir otros tejidos},plural={tumores}}
\newglossaryentry{simulador}{name={simulador de progresión tumoral},description={Programa desitinado a la simulacioón de la evolución de poblaciones de células cancerígenas con reproducción asexual dadas unas condiciones iniciales. Tienen por finalidad el estudio de los tumores y la predicción de los escenarios futuros del tumor},plural={simuladores de progresión tumoral}}
\newglossaryentry{gen}{name={gen},description={Unidad básica de material genético}, plural={genes}}
\newglossaryentry{genotype}{name={genotipo},description={Conjunto de genes que compone la totalidad del material genético de una célula}, plural={genotipos}}
\newglossaryentry{mut}{name={mutación},description={Alteración permanente de alguno de los genes que compone el genotipo de una célula}, plural={mutaciones}}
\newglossaryentry{fitness}{name={fitness},description={Métrica que determina el grado de adaptación de una determinada célula al entorno en que se encuentra, y, por tanto, determina las posibilidades que tiene de sobrevivir, reproducirse o morir modificando así la población de su genotipo. En ocasiones está definida por varios parámetros como las tasas de nacimiento y muerte}, plural={fitnesses}}
\newglossaryentry{tasaNac}{name={tasa de nacimiento},description={Tasa que se emplea en un distribución exponencial para determinar el tiempo en que se producirá el siguiente nacimiento de una célula de este genotipo}, plural={tasas de nacimiento}}
\newglossaryentry{tasaDef}{name={tasa de muerte},description={Tasa que se emplea en un distribución exponencial para determinar el tiempo en que se producirá la siguiente muerte de una célula de este genotipo}, plural= {tasas de muerte}}
\newglossaryentry{tasaMut}{name={tasa de mutación},description={Tasa que se emplea en un distribución exponencial para determinar el tiempo en que se producirá la próxima mutación de una célula de este genotipo}, plural={tasas de mutación}}
\newglossaryentry{epistasis}{name={epistasia},description={Fenómeno por el cual la presencia de un determinado gen mutado afecta a la viabilidad de otro gen, haciendo que el fitness de un gen se vea afectado por la presencia de otro gen mutado en el mismo genotipo}, plural={epistasis}}
\newglossaryentry{tad}{name={TAD},description={"Tipo Abstracto de Datos". Es el conjunto de una estructura en la que almacenan datos, y las funciones implementadas para manipularlos.}, plural={TADs}}

% Rerefencias
\bibliography{src/bibliografia}
\usepackage{graphicx}
\usepackage{listings}

\usepackage{algorithm}
\usepackage{algorithmic}



\usepackage[T1]{fontenc}
\usepackage[scaled=0.84]{beramono}
\usepackage[utf8]{inputenc}

\lstnewenvironment{CPP}
{\lstset{ %
  language=C++,                
  basicstyle=\small,     
  stepnumber=1,                                       
  numbersep=5pt,                 
  backgroundcolor=\color{white},  
  showspaces=false,             
  showstringspaces=false,         
  showtabs=false,                 
  frame=single,                   
  rulecolor=\color{black},       
  tabsize=4,                               
  keywordstyle=\color{blue},          
  commentstyle=\color{brown},    
  stringstyle=\color{red},       
  escapeinside={\%*}{*)},        
  morekeywords={*,...}              
}}{}

\lstnewenvironment{R}
{\lstset{ %
  language=R,                
  basicstyle=\small,     
  stepnumber=1,                                       
  numbersep=5pt,                 
  backgroundcolor=\color{white},  
  showspaces=false,             
  showstringspaces=false,         
  showtabs=false,                 
  frame=single,                   
  rulecolor=\color{black},       
  tabsize=4,                                         
  keywordstyle=\color{blue},          
  commentstyle=\color{brown},    
  stringstyle=\color{red},       
  escapeinside={\%*}{*)},        
  morekeywords={*,...}              
}}{}



% Inicio del documento
\begin{document}

% Elección del idioma (español)
\selectlanguage{spanish}

%
% Portada
%
\pagenumbering{gobble}
\include{portada}
\hypersetup{pageanchor=true}

% Estilo de párrafo de los capítulos
\setlength{\parskip}{0.75em}
\renewcommand{\baselinestretch}{1.25}
% Interlineado simple
\spacing{1}

%
% Agradecimientos
%
\pagenumbering{Roman}
\setcounter{page}{0}
\chapter*{Agradecimientos}

A mi familia y amigos por su apoyo durante estos meses.

A Ramón Díaz Uriarte por su dedicación y guía a lo largo de este proyecto.

% Cita
\begin{flushright}
\textit{``Lo que importa verdaderamente en la vida no son los objetivos que nos marcamos, sino los caminos que seguimos para lograrlo.''}

Peter Bamm
\end{flushright}
  

%
% Resumen
%
% Resumen en inglés
\chapter*{Abstract}

\begin{abstractEn}
The increasing frecuency of cancer cases, as well as its severity, has made the cientific comunity amplify its efforts on researching this illness. Due to this increased interest, numerous tools have appeared to aid with this purpose, some of them, focused on predicting the evolution of cancerous tumours, are known as tumoral progression simulators.

On this project, the functionality of one of these simulators, \textbf{OncoSimulR} is going to be extended. A group of changes is going to be added to the program in order to allow the user to specify arbitrarily complex user variables that depend on existing parametesrs of the simulation, giving the user the possibility of tracking any data of the tumoral progression. This variables will also be usable for the definition of the interventions that the program can already emulate, adapting them to the current state of the tumor in a given moment, in order to achieve the simulation of adaptive therapy.

Furthermore, some tests will be developed in order to ensure the optimal performance of this new functionality. Lastly, examples of some use cases of the program will be given, illustrating the new possible ways to use the tool.
\end{abstractEn}

% Palabras clave en inglés
\begin{keywordsEn}
Cancer, tumor, tumoral progression simulator, OncoSimulR, user variables, interventions, adaptive therapy.
\end{keywordsEn}

% Resumen en español
\chapter*{Resumen}

\begin{abstractEs}
La creciente preocupación por la proliferación del \gls{cancer} como enfermedad, así como por la gravedad de este, ha propiciado que en los últimos años se hayan multiplicado los esfuerzos en su investigación. Con este fin han surgido numerosas herramientas para predecir la evolución de un \gls{cancer}, conocidas cómo \glspl{simulador}.

En este trabajo se va a ampliar la funcionalidad de uno de estos simuladores, \textbf{OncoSimulR}. Sobre éste se va a introducir la posibilidad de que el usuario defina una serie de variables de usuario tan complejas como desee, en función de parámetros de la simulación, que permitirán al usuario obtener prácticamente cualquier dato que quiera acerca de la evolución del \gls{tumor}. Además se va a permitir que estas variables se empleen en la definición de las intervenciones que esta aplicación permite reproducir, con el fin de adaptar estas al estado del \gls{tumor} en cada momento, logrando de este modo emular la conocida como terapia adaptativa.

Tras explicar en qué consisten los cambios llevados a cabo para esto, se describirá la batería de pruebas empleada para comprobar su correcto funcionamiento, y se mostrarán algunos ejemplos prácticos del uso del programa, ilustrando esta nueva funcionalidad.
\end{abstractEs}

% Palabras clave en español
\begin{keywordsEs}
\Gls{cancer}, \gls{tumor}, \gls{simulador}, OncoSimulR, variables de usuario, intervenciones,terapia adaptativa.
\end{keywordsEs}


%
% Glosario
%
\printglossary[title=Glosario,toctitle=Glosario]
\printglossary[title=Acrónimos,toctitle=Acrónimos,type=\acronymtype]

% Estilo de párrafo de los índices
\setlength{\parskip}{1pt}
\renewcommand{\baselinestretch}{1}

%
% Tabla de contenidos
%
\tableofcontents
\listoftables
\listoffigures
\cleardoublepage

% Estilo de párrafo de los capítulos
\setlength{\parskip}{0.75em}
\renewcommand{\baselinestretch}{1.25}
% Interlineado simple
\spacing{1}
% Numeración contenido
\pagenumbering{arabic}
\setcounter{page}{1}

%
% Introducción
%
\chapter{Introducción}

\section{Motivación}
\begin{itemize}
	\item ¿Intro sobre el cáncer e intervenciones?
	\item ¿Terapia adaptativa?
\end{itemize}



\section{Objetivos}

\begin{itemize}
	\item Familiarizarse con OncoSimulR (esopecialmente el flujo general de simulación y las intervenciones)
	\item Aprender sobre R y C++
	\item Completar el merge de las ramas anteriores solucionando todos los errores que surjan
	\item ¿Optimizar el código de intervenciones?
	\item Añadir la posibilidad de definir variables de usuario arbitrariamente complejas tanto como salida en sí mismas como para emplearse en las intervenciones
	\item Generar una batería de pruebas para comprobar el correcto funcionamiento de las variables de usuario y las intervenciones que las emplean
\end{itemize}

\section{Estructura del documento}

TODO: Descripción de la estructura del documento
%
% Estado del arte
%
\chapter{Estado del Arte\label{sec:estado_del_arte}}





%
% Diseño
%
\chapter{Diseño\label{sec:disenho}}

\section{Introducción. Terapia adaptativa y variables de usuario}

Como hemos visto en la sección 1.1 (TODO), la base fundamental de la terapia adaptativa es la posibilidad de crear tratamientos específicos para determinados momentos en el desarrollo de un tumor. Dentro del programa \textbf{OncoSimulR} los tratamientos están simulados a traves de las \textbf{intervenciones}. Por tanto, para simular la terapia adaptativa debemos lograr que estas intervenciones puedan depender del  estado de desarrollo del tumor, es decir, del estado de la simulación en el momento concreto en que se va a ejecutar la intervención.

Para ello se ha decidido añadir al programa la posibilidad de definiri variables de usuario, es decir, se ha incluido la posibilidad de que el usuario defina una serie de variables cuyo valor sea una función arbitrariamente compleja de distintos parámetros de la simulación. A lo largo de la simulación tumoral, estas variables se irán actualizando de acuardo a la fórmula introducida por el usuario, haciendo que los valores vayan variando en función de la evolución del tumor. 

Estas variables podrán luego ser empleadas para elaborar las expresiones que definen la acción a realizar en una intervención (como se ha visto en el punto 2 TODO), permitiendo así condicionar el efecto de una intervención al estado del tumor en el momento de la ejecución de esta, es decir, permitiendo la terapia adaptativa.

Con esto en mente procedemos al diseño de estas variables de usuario.

\section{Definición de variables y reglas}
A la hora de plantearse el diseño necesario para esta funcionalidad, lo primero que necesitamos saber es qué son y como se definen las variables de usuario que vamos a implementar.

Las variables de usuario son variables arbitrariamente complejas que un usuario puede definir a partir de combinaciones o funciones de algunas variables ya existentes en el programa. Para obtener estas variables arbitrariamente complejas vamos a necesitar de una serie de reglas que determinen el valor de estas variables en cada momento de acuerdo a las condiciones propuestas por el usuario para cada una de las variables.

Por tanto podemos determinar que por un lado tendremos un listado de variables, cada una de las cuales debera contener:
\begin{itemize}
	\item \textbf{Nombre}: nombre de la variable de usuario que permite identificarla y diferenciarla de el resto de variables definidas.
	\item \textbf{Valor}: valor numérico de la variable en cada momento.
\end{itemize}

Por otro lado tendremos un listado con las reglas que permiten modificar y actualizar las distintas variables definidas. Estas reglas contan de:

\begin{itemize}
	\item \textbf{ID}: identificador único de cada regla que permite diferenciarla del resto de reglas definidas.		
	\item \textbf{Condition}: la condición que, si se cumple, hará que la variable de usuario se modifique.
	\item \textbf{Action}: acción que se llevará a cabo en caso de que la \textbf{Condition} se cumpla, es decir, es el conjunto de modificaciones a realizar a las variables de usuario en caso de que se cumpla la condición de la regla. podrá contener una o varias fórmulas que determinarán los nuevos valores de una o varias variables de usuario.
\end{itemize}

En este punto debemos mencionar también, qué variables del programa se ha decidido que sean empleables para la modificación de variables y definición de las reglas por parte del usuario y los motivos por los que se ha decidido esto. Las variables disponibles para la definición de las reglas son:

\begin{itemize}
	\item \textbf{birthRate}: la tasa de nacimiento correspondiente a lcada genotipo, es decir,es un nñumero que determina la probabilidad que tiene una célula con este genotipo de reproducirse, y por tanto aumentar su número total en 1 en un ciclo de tiempo.
	\item \textbf{deathRate}: la tasa de muerte correspondiente a cada genotipo, es decir, es un número que determina la probabilidad que tiene una célula con este genotipo de morir, y por tanto reducir su número total en 1 en un ciclo de tiempo.
	\item \textbf{mutationRate}: la tasa de mutación correspondiente a cada genotipo, es decir, es un número que determina la probabilidad que tiene una célula con este genotipo de mutar, y por tanto variar alguno de sus genes, pasando de ser de este genotipo a ser de otro genotipo distinto en 1 en un ciclo de tiempo.
	\item \textbf{N}: número total de células en la simulación, teniendo en cuenta todos los genotipos existentes.
	\item \textbf{n\_x}: número de células de un determinado genotipo concreto, por ejemplo, \textbf{n\_A} será el número de células del genotipo A, y \textbf{n\_A\_B} las del genotipo A\_B
	\item \textbf{T}: tiempo transcurrido desde el inicio de la simulación, en segundos, con precisión de milisegundos.
	\item \textbf{Variables de usuario}: otras variables de usuario definidas en la misma simulación.		
\end{itemize}

Estas variables se han añadido con la finalidad de dar a las variables de usuario la mayor flexibilidad posible para que el usuario las defina de acuerdo a sus necesidades, considerando todas las variables interesantes tanto para la terapia adaptativa como desde un punto de vista biológico para estudiar su evolución.

Las variables \textbf{N}, \textbf{n\_x} y \textbf{T} se consideran necesarias para la terapia adaptativa, ya que permiten definir variables de usuario que se modifiquen de auerdo a la evolución del tamaño del tumor (\textbf{N}), a la cantidad de las células de un determinado tipo, que se consideren más perjudiciales o sensibles a un determinado tratamiento (\textbf{n\_x}) o al tiempo transcurrido (\textbf{T}). Estas variables definidas de acuerdo a estos atributos se podrían emplear luego para definir las intervenciones que permite el programa OncoSimulR, permitiendo adaptarlas de forma adecuada a las distintas combinaciones posibles de estos 3 atributos a lo largo de la simulación, generando así la posibilidad de simular terapia adaptativa. A estas 3 variables se les suman las propias variables de usuario, para poder así generar combinaciones aún más complejas que permitan simular escenarios altamente enrevesados.

Las otras variables  \textbf{birthRate}, \textbf{deathRate} y \textbf{mutationRate}, se han añadido con una finalidad más informativa, ya que, a pesar de no ser tan interesantes desde el punto de vista de la terapia adaptativa (aunque pueden aplicarse en algunos casos), dan una información importante acerca de la adaptabilidad de los distintos genotipos a lo largo de la simulación, haciendo que el estudio de la evolución de variables de usuario que sean combinación de estas variables a lo largo de la simulación pueda tener un gran interés para el estudio tumoral.

\section{Algoritmo de actualización}

El algoritmo de actualización de las variables de usuario es a priori, bastante sencillo, pero su principal complejidad radica en su implementación e integración en el flujo de código principal, como veremos en la sección 4.

En rasgos generales el algoritmo de actualización es:

\begin{enumerate}
	\item El usuario define las varibales de usuario que desea y las reglas para modificarlas
	\item Se verifica que estan definidas correctamente
	\item Cada unidad de tiempo muestreada en la simulación:
	\begin{enumerate}
		\item Para cada regla:
		\begin{enumerate}
			\item Si se cumple la Condition:
				\begin{enumerate}
					\item Modificamos el valor de las variables de usuario según lo especificado por el usuario en el atributo Action		
				\end{enumerate}
			\item Se pasa a la siguiente regla volviendo a 1)	
		\end{enumerate}
		\item Se procede a la ejecución de las intervenciones con los valores de las variables calculados	
	\end{enumerate}		
\end{enumerate}

\section{Definición de la estructura de salida}

Para obtener los datos de evolución de las variables de usuario definidas en el tiempo a lo largo de la simulación, se deberán modificar la salida del programa de simulación para que incluya estos. Se debe añadir por tanto a la salida
\begin{itemize}
	\item Una lista con los nombres (``ids'') de las variables de usuario creadas.
	\item Una matriz con los valores de las variables de usuario en cada unidad de tiempo muestreada (en el mismo orden que los nombres de las variables en en la lista de nombres de las variables de usuario).	
\end{itemize}

La finalidad de añadir estos datos a la salida de la simulación es poder emplear las variables de usuario (que son arbitrariamente complejas) como datos de salida y no solo como valores a emplear en otros procesos como las intervenciones. Por ejemplo, obtener estos datos permite al usuario graficar la evolución en el tiempo de algunas variables de usuario creadas, lo que puede ser de gran interés para evaluar la evolución del tumor en la simulación. Además, la adición de estos datos en la salida facilitará enormemente las labores de testeo de la funcionalidad implementada.	





%
% Desarrollo
%
\chapter{Desarrollo\label{sec:desarrollo}}

\section{Introducción}
TODO
¿Ramas?
¿Exprtk?
¿Rcpp?
¿R y C++?

\section{Unificación de Ramas}

Cómo se ha explicado en el punto 2 (TODO) al iniciar este trabajo se partía de 2 ramas con versiones distintas del programa OncoSimulR cada una con una funcionalidad recientemente implementada en cada una de ellas. Estas 2 ramas debían unificarse para proceder a la implementación de la nueva funcionalidad sobre la versión combinada de ambas.

La unificación de ambas ramas se realizó de manrea manual, ya que el merge automático proporcionado por la herramienta Git generaba un gran número de conflictos y no respetaba la funcionalidad de las ramas. Al proceder, se decidió realizar el merge sobre la rama que contenía la funcionalidad del birth y death rate (explicada en el punto 2 TODO) por ser la que más modificaciones introducía sobre el código base de la simulación, y sobre esta se introdujeron los cambios correspondientes a la rama con la funcionalidad de las intervenciones (explicada en el punto 2 TODO).

Para realizar la unificación de ambas ramas se procedió del siguiente modo:
\begin{enumerate}
	\item Sobre la rama que se ha indicado anteriormente se añadieron los ficheros de código R y C++ exclusivos de la funcionalidad de intervenciones (intervention.cpp, intervention.h, intervention.R, multivarioant\_hypergeometric.cpp y multivariant\_hypergeometric.h), los ficheros de testeo de esta funcionalidad (test.intervention.R y test.Z-intervention.R) y los ficheros de documentación y viñetas que expican e ilustran el funcionamiento de esta.
	\item Tras esto se procedió a la modificación de la función updatePopulations del fichero intervention.cpp, encargada de actualizar las poblaciones de los distintos genotipos tras la ejecución de una intervención, para que emplease el birth y death rate implementado en la funcionalidad de la otra rama.
	\item Se actualizó el flujo de código principal del algoritmo BNB de Mahler (TODO: explicado en el punto 2) en el fichero BNB\_nr.cpp para añadir las llamadas necesarias para añadir la funcionalidad correspondiente a la ejecución de las intervenciones.
	\item Por último se modificaron las cabeceras y llamadas a todas las funciones del programa OncoSimulR para que tomasen de forma adecuada el argumento interventions, teniendo especial cuidado al tratar con los ficheros RcppExports.R, RcppExports.cpp y OncoSimulR\_init.c, encargados de la comunicación entre los ficheros en R y en C++ mediante la librería Rcpp.
\end{enumerate}

Posteriormente se realizó un merge automático con una tercera rama modificada por el tutor de este trabajo para incluir unas pequeñas modificaciones en la documentación del programa.

Al finalizar este proceso, se procedió a la comprobación de que la funcionalidad de ambas ramas era correcta en la rama combinada mediante la ejecución de los tests pertinentes, y a asegurar que el merge no ha afectado al resto de características del programa. Todo este testeo se explicará de forma más detallada en el punto 5.1 TODO.

\section{Optimización del código de intervenciones}
Durante la realización del merge de las ramas previas, se encontraron algunos problemas de optimización en el código de las intervenciones que hacían que el tiempo de ejecución fuese más alto de lo deseado, por tanto antes de proseguir con la nueva funcionalidad se procedió a solucionar esto.

El principal problema detectado en este aspecto eran fragments de código que se ejecutaban en todas las iteraciones, independeintemente de si se realizaban intervenciones o no, cuando solo era necesaria su ejecución cuando estas si se realizaban. Si bien estos fragmentos de código no eran en sí mismos muy costosos coimputacionalmente, en una simulación de gran tamaño, con un gran número de iteraciones (como las que se realizan habitualmente en est programa), estos fragmentos se estarían ejecutando de forma innecesaria un enorme número de veces, haciendo que la variación en el tiempo de ejecución sí fuese bastante significativo.

Los cambios realizados, por tanto, fueron:
\begin{itemize}
	\item Modificación de la función executeInterventions de interventions.cpp:
	\begin{itemize}
		\item Cambiamos la función para que devuelva true si ha habido una intervención y false en caso contrario en lugar de devolver void como hasta ahora. Esto nos permite desde el algoritmo principal saber si ha habido o no intervención para ejecutar los fragmentos de código necesarios, además nos permite desde fuera obtener los tiempos en que se ejecutan las intervenciones, facilitándonos el testeo posteriormente.
		\item  Hacemos que en la función sólamente se llame a updatePopulations (que actualiza las poblaciones de cada genotipo) en los casos en los que se ha ejecutado alguna intervención, puesto que son los ñunicos casos en que las poblaciones pueden haber cambiado.
	\end{itemize}
	\item Modificamos también la función ParseWhatHappens para que devuelva void en vez de true/false. Antes devolvía true si se ejecutaba sin errores y en caso de error lanzaba una excepción y devolvía false, lo cual no tenía sentido, por tanto ahora no devuelve nada y en caso de error simplemente lanza una excepción.
\end{itemize}

Tras esto de nuevo procedemos a comprobar que la funcionalidad es correcta mediante los tests explicados en el punto 5 TODO, y procedemos a comenzar con la nueva funcionalidad.

\section{Implementación de variables de usuario y reglas}

\subsection{En R}
Desde R debemos hacer posible que el usuario pueda definir las variables de usuario que desee, y las reglas que quiere que se emplee para su cálculo y actualización a lo largo de la simulación, así como permitir que pase estas como argumento a la funcion OncoSimulIndiv de OncoSimulR, que es la función responsable de llevar a cambo la simulación.

Para ello en primer lugar vemos como se deben definir las variables de usuario y las reglas en R:

\begin{itemize}
	\item Variables de usuario: las variables de usuario hemos visto en el punto 3.1 que simplemente contienen un nombre que identifica a la variable, y un valor numérico. Por tanto el usuario definirá en R una simple lista con todas las variables de usuario que desee crear para la simulación con estos atributos, del siguente modo:

	\begin{itemize}
		\item Name: string que contiene el nombre que identificará a la variable, debe ser único para la variable.
		\item Value: valor inicial que tomará la veriable al inicio de la simulación.
	\end{itemize}

	\item Reglas: del mismo modo hemos visto en el 3.1 que las reglas constan de un ID, un Condition y una Action. De igual forma el usuario definirá estas reglas ediante una lista en R de la forma:

	\begin{itemize}
		\item id: string identifica la regla y la distingue del resto, debe ser único para cada regla.
		\item Condition: expresión con valor true/false que determina si se ejecuta o no la regla. Puede emplear las variables especificadas en el punto 3.1.
		\item Action: una o varias expresiones con valor numérico separadas por ";" que determinan el valor que tomarán las variables de usuario en caso de que se ejecute la rega. Del mismo modo que el campo Condition, puede emplear las variables especificadas en el punto 3.1.
	\end{itemize}
\end{itemize}

Partiendo de estas definiciones, creamos en R una serie de funciones que comprueban que el usuario ha especificado las variables de usuario y reglas de forma correcta y las adaptan para que puedan pasartse a C++ de forma correcta mediante la librería Rcpp. Estas funciones son en el caso de las variables de usuario:

\begin{itemize}
	\item check\_same\_name: comprueba que en la lista definida no existen varias variables de usuario con el mismo nombre.
	\item verify\_user\_vars: comprueba que todas las variables de usuario definidas en la lista tienen todos los atributos necesarios (Name y Value).
	\item createUserVars: mediante llamadas a las 2 funciones anteriores comprueba que la lista de variables de usuario definida es correcta, y si lo es, la devuelve.
\end{itemize}

Y en el caso de las reglas:

\begin{itemize}
	\item check\_double\_rule\_id: comprueba que en la lista de reglas definida no hay 2 reglas con el mismo id.
	\item check\_action: comprueba la correcta especificación de la expresión del campo Action, para ello comprueba que, si hay varias expresiones, estas están correctamente separadas por ";" y que cada una de las expresiones tiene la forma: <variable de usuario> = <expresión>.
	\item verify\_rules: comprueba que las reglas contienen todos los atributos necesarios (id, Action y Condition), y que estos son correctos llamadno a las 2 funciones anteriores. 
	\item transform\_rule: modifica las expresiones Action y Condition de las reglas, ya que aunque en estas las variables n\_X (explicadas en 3.1) se especifican como n\_A, n\_B, etc; en C++ los genotipos se definen con nñumeros y no con letras, por tanto esta función transforma n\_A, n\_B... en n\_1, n\_2... 
	\item adapt\_rules\_to\_cpp: mediante llamadas a la función anterior, modifica las expresiones action y condition de las distintas reglas definidas en la lista para adaptarlas a C++.
	\item createRules: mediante llamadas a verify\_rules y a adapt\_rules\_to\_cpp comprueba que las reglas definidas son correctas y las adapta a C++ devolviendo la lista de reglas ya transformada.
\end{itemize}

Por último, en R, se deben modificar las llamadas y la cabecera de la función OncoSimulIndiv y de las demás funciones internas de OncoSimulR para que reciban de forma adecuada los nuevos argumentos userVars (lista de variables de usuario) y rules (listado de reglas).

\subsection{En C++}

\begin{itemize}
	\item estructuras Rule y UserVarsInfo
	\item funciones de user\_var.cpp:
	\begin{itemize}
		\item Básicas (createRule, createUserVarsInfo, destroyRule, compareRules, addRule, printRule, printUserVarsInfo)
		\item isValid (Comprueba si las expresiones de las acciones de las reglas están bien formadas)
		\item executeRules (Comprueba si se cumple la condición mediante exprtk, y si se cumple llama a parseAction)
		\item parseAction ("compila" la acción de la regla con exprtk y la ejecuta variando las variables de usuario necesarias en el map de UserVarsInfo)
	\end{itemize}
\end{itemize}

\section{Integración en el flujo principal de código}

\begin{itemize}
	\item Se incluye la ejecución de las reglas en el flujo de código de la simulación en la función nr\_innerBNB de BNB\_nr.cpp justo antes de la ejecución de las intervenciones para permitir que las variables de usuario estén actualizadas antes de realizar las intervenciones ya que se emplean para estas.
	\item Se crea la función evalFVarsFitness que permite acceder al birthrate, deathrate y mutationrate que se deben pasar como argumentos  para executeRules para poder añadirse a la tabla de símbolos de exprtk y poder emplearse en las fórmulas para las codiciones o modificaciones de las reglas para actualizar las variables de usuario.
	\item Empleamos tambien la función evalFVars para acceder a las poblaciones de los genotipos con la misma finalidad.
	\item Modificamos la estructura de salida de la función BNB\_nr\_Algo5 para que incluya la lista con los nombres de las variables de usuario y la matriz con los valores de estas en cada sampleEvery
	\item Modificamos las cabeceras y llamadas a funciones para que incluyan las nuevas variables userVars, Rules y userVarsInfo, y se realicen las llamadas de forma adecuada.
	\item Modificación de los parametros de entrada de todas las funciones de OncoSimulR para que tomen de forma adecuada el parámetro de intervenciones, así como actualización de RcppExports.R, RcppExports.cpp y OncoSimulR\_init.c con el mismo fin.
\end{itemize}

\section{Inclusión de las variables en las intervenciones}

\begin{itemize}
	\item Modificamos executeInterventions para que reciba como argumento la estructura userVarsInfo con los datos de las variables de usuario y añadimos un pequeño fragmento que añade las variables de usuario y sus valores a la tabla de símbolos de exprtk para poder ser empleadas en las imtervenciones.
	\item Realizamos el mismo proceso para ParseWhatHappens.
	\item Modificamos las cabeceras y llamadas a estas 2 funciones de forma adecuada para que incluyan userVarsInfo en todas las llamadas. 
\end{itemize}



%
% Resultados
%
\chapter{Pruebas y Resultados\label{sec:resultados}}

\section{Comprobación de funcionalidad previa}

Antes de proceder a introducir toda la nueva funcinalidad explicada a partir del punto 4.4, y tras realizar la unificación de las 2 ramas de las que se partía y la solución de los problemas de optimización explicados en 4.3 ,aseguramos que la funcionalidad anterior no se ha visto afectada mediante el uso de la batería de tests ya disponibles para el programa \textbf{OncoSimulR}. Esta batería de tests incluye por una lado unos tests manuales muy extensos y costosos computacionalmente, y otro grupo de tests que comprueban la funcionalidad básica del programa mediante el paquete \textbf{testthat} del lenguaje R. Al realizar esto se detectaron una serie de errores previos al merge en una de las ramas. Los errores detectados están localizados en los tests de la funcionalidad de intervenciones, y son errores en los propios tests y no fallos de funcionalidad. Para solucionar estos errores:

\begin{itemize}
		\item Se modifica la salida del algoritmo de simulación BNB implementado en la función \textbf{nr\_BNB\_Algo5} para que incluya una lista con los tiempos de la simulación en los que se ejecutan intervenciones (esto nos permite buscar los puntos exactos donde se realizan las intervenciones en los tests). De este modo la salida de la función principal de simulación \textbf{OncoSimulIndiv} ahora incluye el campo \textbf{x\$other\$interventionTimes} con estos tiempos (donde ``x'' es la variable donde se almacena la salida de la función).
		\item Adelantamos el punto de ejecución de las intervenciones en flujo principal del código de la simulación realicen antes de rellenar las estructuras de salida. Con esto conseguimos que al actualizarse las poblaciones y tasas de nacimiento y muerte al final de cada ciclo de simulación, estas incluyan los cambios realizados por las intervenciones; antes de esta modificación, los cambios que realizaban las intervenciones no se veían reflejados en las poblaciones mostradas a la salida hasta el tiempo correspondiente a la siguiente iteración de la simulación, lo que generaba un desfase entre lo que se testeaba (que emplea los datos de salida) y las poblacioes internas empleadas en la simulación (que sí incluían los cambios realizados por las intervenciones).
		\item Comentar la comprobación de totPopSize != totalPop en reduceTotalPopulation (interventions.cpp) porque al haberlo adelantado totPopSize no está actualizado.
		\item Añadimos un fragmento de código para eliminar los genotipos con población 0 despues de realizar intervenciones. Esto estaba generando un error pues algunas operaciones de la simulación no se pueden realizar con poblaciones 0. 
		\item Modificamos los tests ya existentes de intervenciones para adaptarlos a estos nuevos cambios y comprobamos que todos los tests se pasan correctamente. Estas modificaciones consisten principalmente en hacer que se empleen los tiempos almacenados en \textbf{x\$other\$interventionTimes} para buscar los puntos en los que se realizan intervenciones y comprobar que estas son correctas, en lugar de buscarlas recorriendo la estructura de salida con las poblaciones en cada unidad de tiempo, pues este método no permitía un testeo correcto además de ser mucho menos eficiente.
\end{itemize}

Por último, tras implementar toda la nueva funcionalidad (explicada a partir del punto 4.4) volvemos a ejecutar toda la batería de tests  para asegurar, de nuevo, que estos cambios no han afectado al correcto funcionamiento de la funcionalidad previa.

\section{Tests de nueva funcionalidad}

Para probar el correcto funcionamiento de la nueva funcionalidad implementada que incluye la creación de variables de usuario y su actualización, así como la simulación de la terapia adaptativa mediante la inclusión de estas variables en las intervenciones; se ha ampliado la batería de tests del programa (concretamente los tests que emplean la herramienta \textbf{textthat} incluyendo un nuevo documento de tests \textbf{test.user\_var.R} que prueba la funcionalidad de las variables de usuario, y se ha ampliado el documento de tests \textbf{test.intervention.R} que incluye los tests sobre intervenciones, para añadir los casos en que estas dependen de los valores contenidos en variables de usuario. Ambos ficheros están incluidos en el directorio \textbf{tests/testthat}.

\subsection{Variables de usuario}

Como hemos comentado los tests que comprueban el correcto funcionamiento de las variables de usuario están contenidos en el fichero \textbf{test.user\_var.R}, estos tests contienen las siguientes comprobaciones.
\begin{itemize}
	\item Las variables de usuario se crean correctamente
	\item 2 variables no pueden tener el mismo nombre
	\item Las reglas para la actualización de las variables se crean correctamente
	\item 2 reglas no pueden tener el mismo ID
	\item Al ejecutarse las reglas, las variables de usuario se modifican de forma correcta (en función de la variable T)
	\item Al ejecutarse las reglas las variables de usuario se modifican de forma correcta (en función de la variable N)
	\item Al ejecutarse las reglas las variables de usuario se modifican de forma correcta (en función de las poblaciones por genotipo n\_x)
	\item Al ejecutarse las reglas las variables de usuario se modifican de forma correcta (en función de la otras variables de usuario)
\end{itemize}

\subsection{Intervenciones con variables de usuario}

Una vez hemos comprobado que las variables de usuario funcionan de forma adecuada, ampliamos los tests del fichero \textbf{test.intervention.R} para asegurar que las intervenciones funcionan de forma correcta cuando depende de las variables de usuario creadas. Para ello añadimos 2 nuevos tests en este fichero para comprobar:
\begin{itemize}
	\item Las intervenciones se ejecutan de forma correcta cuando el \textbf{Trigger} (condición de ejecución de la intervención) depende de variables de usuario.
	\item Las intervenciones se ejecutan de forma correcta cuando el \textbf{WhatHappens} (modificación en poblaciones o tasas realizada por la intervención) depende de variables de usuario.
\end{itemize}

\section{Resultados}

Tras todo el trabajo realizado, en esta sección vamos a comprobar que, efectivamente, los resultados obtenidos cumplen con lo esperado e, incluso, se han logrado resultados positivos que no se habían planteado incialmente como objetivos del trabajo.

\subsection{Terapia Adaptativa}
Como se ha comentado a lo largo del documento, el principal objetivo de todas las modificaciones implementadas es lograr la simulación de la terapia adaptativa, explicada en detalle en el punto 2 (TODO). 
\subsection{Variables de usuario como datos de interés a la salida}
\subsection{Optimización en la ejecución de intervenciones}

%
% Conclusiones
%
\chapter{Conclusiones y trabajo futuro\label{sec:conclusiones}}

\section{Conclusiones}

\section{Líneas de trabajo futuro}

\begin{itemize}
	\item Emplear variables de usuario en otras funcionalidades más allá de las intervenciones
	\item Permitir graficar variables de usuario en función del tiempo como una salida interesante de forma similar a la evolución de la población de los genotipos
\end{itemize} 

%
% Página en blanco
%
\cleardoublepage

%
% Bibliografía
%
\printbibliography[heading=bibintoc]

% No expandir elementos para llenar toda la página
\raggedbottom

%
% Apéndices
%
\appendix
\cleardoublepage
\addappheadtotoc
\appendixpage

%
% TODO: Apéndices del TFG
%
\chapter{Código para  reproducir los ejemplos de la sección Resultados \label{sec:ejemplos}}
\section{Ejemplo  de uso de variables de usuario para controlar la proporcion de cada genotipo en la población total}
Código para reproducir el segundo ejemplo del punto 5.3.1:
\begin{R}
dfuv2 <- data.frame(Genotype = c("WT", "B", "A", "B, A", "C, A"),
                       Fitness = c("0*n_",
                                    "1.5",
                                    "1.002",
                                    "1.003",
                                    "1.004"))

afuv2 <- allFitnessEffects(genotFitness = dfuv2,
                          frequencyDependentFitness = TRUE,
                          frequencyType = "abs")

userVars <- list(
    list(Name           = "genAProp",
        Value       = 0.5
    ),list(Name           = "genBProp",
        Value       = 0.5
    ),list(Name           = "genABProp",
        Value       = 0.0
    ),list(Name           = "genACProp",
        Value       = 0.0
    )
)

userVars <- createUserVars(userVars)

rules <- list(
        list(ID = "rule_1",
            Condition = "TRUE",
            Action = "genBProp = n_B/N"
        ),list(ID = "rule_2",
            Condition = "TRUE",
            Action = "genAProp = n_A/N"
        ),list(ID = "rule_3",
            Condition = "TRUE",
            Action = "genABProp = n_A_B/N"
        ),list(ID = "rule_4",
            Condition = "TRUE",
            Action = "genACProp = n_A_C/N"
        )
    )

rules <- createRules(rules, afuv2)

set.seed(1)
uvex2 <- oncoSimulIndiv(
                    afuv2, 
                    model = "McFLD",
                    mu = 1e-4,
                    sampleEvery = 0.01,
                    initSize = c(20000, 20000),
                    initMutant = c("A", "B"),
                    finalTime = 10,
                    onlyCancer = FALSE,
                    userVars = userVars,
                    rules = rules
                    )

plot(uvex2, show = "genotypes", type = "line")      

plot(unlist(uvex2$other$userVarValues) [c(FALSE, FALSE, FALSE, FALSE, TRUE)], unlist(uvex2$other$userVarValues) [c(TRUE, FALSE, FALSE, FALSE, FALSE)], xlab="Time", ylab="Proportion", ylim=c(0,1), type="l", col="purple")  
lines(unlist(uvex2$other$userVarValues) [c(FALSE, FALSE, FALSE, FALSE, TRUE)], unlist(uvex2$other$userVarValues) [c(FALSE, TRUE, FALSE, FALSE, FALSE)], type="l", col="#E6AB02") 
lines(unlist(uvex2$other$userVarValues) [c(FALSE, FALSE, FALSE, FALSE, TRUE)], unlist(uvex2$other$userVarValues) [c(FALSE, FALSE, TRUE, FALSE, FALSE)], type="l", col="#1B9E77") 
lines(unlist(uvex2$other$userVarValues) [c(FALSE, FALSE, FALSE, FALSE, TRUE)], unlist(uvex2$other$userVarValues) [c(FALSE, FALSE, FALSE, TRUE, FALSE)], type="l", col="#666666") 
legend(0,1,legend=c("genABProp", "genACProp", "genAProp", "genBProp"), col=c("purple", "#E6AB02", "#1B9E77", "#666666"), lty= 1:2) 
\end{R}
\section{Ejemplo  de uso de variables de usuario para calcular la diferencia entre tasa de nacimiento y muerte de los genotipos}
Código para reproducir el segundo ejemplo del punto 5.3.1:
\begin{R}
dfuv3 <- data.frame(Genotype = c("WT", "A", "B"), 
                   Fitness = c("1",
                               "1 + 0.2 * (n_B > 10)",
                               ".9 + 0.4 * (n_A > 10)"
                               ))
afuv3 <- allFitnessEffects(genotFitness = dfuv3, 
                         frequencyDependentFitness = TRUE)

userVars <- list(
    list(Name           = "genWTRateDiff",
        Value       = 0.5
    ),list(Name           = "genARateDiff",
        Value       = 0.5
    ),list(Name           = "genBRateDiff",
        Value       = 0.0
    )
)

userVars <- createUserVars(userVars)

rules <- list(
        list(ID = "rule_1",
            Condition = "TRUE",
            Action = "genWTRateDiff = b_-d_"
        ),list(ID = "rule_2",
            Condition = "TRUE",
            Action = "genARateDiff = b_1-d_1"
        ),list(ID = "rule_3",
            Condition = "TRUE",
            Action = "genBRateDiff = b_2-d_2"
        )
    )

rules <- createRules(rules, afuv3)

set.seed(1)

uvex3 <- oncoSimulIndiv(afuv3,
                       model = "McFLD", 
                       onlyCancer = FALSE, 
                       finalTime = 150,
                       mu = 1e-4,
                       initSize = 5000, 
                       keepPhylog = FALSE,
                       seed = NULL, 
                       errorHitMaxTries = FALSE, 
                       errorHitWallTime = FALSE,
                       userVars = userVars,
                       rules = rules)

plot(uvex3, show = "genotypes", type = "line")    

plot(unlist(uvex3$other$userVarValues) [c(FALSE, FALSE, FALSE, TRUE)], unlist(uvex3$other$userVarValues) [c(FALSE, FALSE, TRUE, FALSE)], xlab="Time", ylab="Rate Diff", xlim=c(0,150), ylim=c(-0.75,0.75), type="l", col="#1B9E77")  
lines(unlist(uvex3$other$userVarValues) [c(FALSE, FALSE, FALSE, TRUE)], unlist(uvex3$other$userVarValues) [c(TRUE, FALSE, FALSE, FALSE)], type="l", col="#A6761D") 
lines(unlist(uvex3$other$userVarValues) [c(FALSE, FALSE, FALSE, TRUE)], unlist(uvex3$other$userVarValues) [c(FALSE, TRUE, FALSE, FALSE)], type="l", col="#666666") 
legend(0,0.75,legend=c("genWTRateDiff", "genARateDiff", "genBRateDiff"), col=c("#1B9E77", "#A6761D", "#666666"), lty= 1:2) 
\end{R}
\section{Ejemplo de terapia adaptativa para atacar al genotipo más adaptado}
Código para reproducir el primer ejemplo del punto 5.3.2:
\begin{R}
dfat2 <- data.frame(Genotype = c("WT", "A", "B"), 
                   Fitness = c("1",
                               "1 + 0.2 * (n_B > 10)",
                               ".9 + 0.4 * (n_A > 10)"
                               ))
afat2 <- allFitnessEffects(genotFitness = dfat2, 
                         frequencyDependentFitness = TRUE)
    

userVars <- list(
    list(Name           = "genADiff",
        Value       = 0
    ),list(Name           = "genBDiff",
        Value       = 0
    ),list(Name           = "genWTDiff",
        Value       = 0
    )
)

userVars <- createUserVars(userVars)

rules <- list(
        list(ID = "rule_1",
            Condition = "TRUE",
            Action = "genADiff = b_1-d_1"
        ),list(ID = "rule_2",
            Condition = "TRUE",
            Action = "genBDiff = b_2-d_2"
        ),list(ID = "rule_3",
            Condition = "TRUE",
            Action = "genWTDiff = b_-d_"
        ),list(ID = "rule_4",
            Condition = "n_A < 1000",
            Action = "genADiff = -1"
        ),list(ID = "rule_5",
            Condition = "n_B < 1000",
            Action = "genBDiff = -1"
        ),list(ID = "rule_6",
            Condition = "n_ < 1000",
            Action = "genWTDiff = -1"
        )
    )

rules <- createRules(rules, afat2)

interventions <- list(
    list(ID           = "i1",
        Trigger       = "genADiff > genBDiff and genADiff > genWTDiff",
        WhatHappens   = "n_A = n_A*0.5",
        Periodicity   = 10,
        Repetitions   = Inf
    ),list(ID           = "i2",
        Trigger       = "genBDiff > genADiff and genBDiff > genWTDiff",
        WhatHappens   = "n_B = n_B*0.5",
        Periodicity   = 10,
        Repetitions   = Inf
    ),list(ID           = "i3",
        Trigger       = "genWTDiff > genADiff and genWTDiff > genBDiff",
        WhatHappens   = "n_ = n_*0.5",
        Periodicity   = 10,
        Repetitions   = Inf
    )
)

interventions <- createInterventions(interventions, afat2)

set.seed(1) ## for reproducibility
atex2 <- oncoSimulIndiv(afat2,
                       model = "McFLD", 
                       onlyCancer = FALSE, 
                       finalTime = 200,
                       mu = 1e-4,
                       initSize = 5000, 
                       keepPhylog = FALSE,
                       seed = NULL, 
                       errorHitMaxTries = FALSE, 
                       errorHitWallTime = FALSE,
                       userVars = userVars,
                       rules = rules,
                       interventions = interventions)

plot(atex2, show = "genotypes", type = "line")

plot(unlist(atex2$other$userVarValues) [c(FALSE, FALSE, FALSE, TRUE)], unlist(atex2$other$userVarValues) [c(TRUE, FALSE, FALSE, FALSE)], xlab="Time", ylab="RateDiff", ylim=c(-0.75,0.75), type="l", col="#A6761D")  
lines(unlist(atex2$other$userVarValues) [c(FALSE, FALSE, FALSE, TRUE)], unlist(atex2$other$userVarValues) [c(FALSE, TRUE, FALSE, FALSE)], type="l", col="#666666") 
lines(unlist(atex2$other$userVarValues) [c(FALSE, FALSE, FALSE, TRUE)], unlist(atex2$other$userVarValues) [c(FALSE, FALSE, TRUE, FALSE)], type="l", col="#1B9E77") 
legend(0,0.75,legend=c("genADiff", "genBDiff", "genWTDiff"), col=c("#A6761D", "#666666", "#1B9E77"), lty= 1:2) 
\end{R}
\section{Ejemplo de terapia adaptativa en escenario verosímil}
Código para reproducir el segundo ejemplo del punto 5.3.2:
\begin{R}
dfat3 <- data.frame(Genotype = c("WT", "A", "B"), 
                   Fitness = c("1",
                               "0.8 + 0.2 * (n_B > 10) + 0.1 (n_A > 10)",
                               "0.8 + 0.25 * (n_B > 10)"
                               ))
afat3 <- allFitnessEffects(genotFitness = dfat3, 
                         frequencyDependentFitness = TRUE)
    

userVars <- list(
    list(Name           = "lastMeasuredA",
        Value       = 0
    ),list(Name           = "lastMeasuredB",
        Value       = 0
    ),list(Name           = "previousA",
        Value       = 0
    ),list(Name           = "previousB",
        Value       = 0
    ),list(Name           = "lastTime",
        Value       = 0
    ),list(Name           = "measure",
        Value       = 0
    ),list(Name           = "treatment",
        Value       = 0
    )
)

userVars <- createUserVars(userVars)

rules <- list(
       list(ID = "rule_1",
            Condition = "T - lastTime < 10",
            Action = "measure = 0"
        ),list(ID = "rule_2",
            Condition = "T - lastTime >= 10",
            Action = "measure = 1;lastTime = T"
        ),list(ID = "rule_3",
            Condition = "measure == 1",
            Action = "previousA = lastMeasuredA;previousB = lastMeasuredB;lastMeasuredA = n_A;lastMeasuredB = n_B"
        ), list(ID = "rule_4",
            Condition = "TRUE",
            Action = "treatment = 0"
        ), list(ID = "rule_5",
            Condition = "lastMeasuredA + lastMeasuredB > 100",
            Action = "treatment = 1"
        ), list(ID = "rule_6",
            Condition = "lastMeasuredA - PreviousA > 500",
            Action = "treatment = 2"
        ),list(ID = "rule_7",
            Condition = "lastMeasuredB - PreviousB > 500",
            Action = "treatment = 3"
        ), list(ID = "rule_8",
            Condition = "lastMeasuredA - PreviousA > 500 and lastMeasuredB - PreviousB > 500",
            Action = "treatment = 4"
        )
    )

rules <- createRules(rules, afat3)


interventions <- list(
    list(ID           = "basicTreatment",
        Trigger       = "treatment == 1",
        WhatHappens   = "N = 0.8*N",
        Periodicity   = 10,
        Repetitions   = Inf
    ),list(ID           = "treatmentOverA",
        Trigger       = "treatment == 2 or treatment == 4",
        WhatHappens   = "n_B = n_B*0.3",
        Periodicity   = 20,
        Repetitions   = Inf
    ),list(ID           = "treatmentOverB",
        Trigger       = "treatment == 3 or treatment == 4",
        WhatHappens   = "n_B = n_B*0.3",
        Periodicity   = 20,
        Repetitions   = Inf
    ),list(ID           = "intervention",
        Trigger       = "lastMeasuredA+lastMeasuredB > 5000",
        WhatHappens   = "N = 0.1*N",
        Periodicity   = 70,
        Repetitions   = Inf
    )
)

interventions <- createInterventions(interventions, afat3)

set.seed(1) ## for reproducibility
atex2 <- oncoSimulIndiv(afat3,
                       model = "McFLD", 
                       onlyCancer = FALSE, 
                       finalTime = 200,
                       mu = 1e-4,
                       initSize = 5000, 
                       keepPhylog = FALSE,
                       seed = NULL, 
                       errorHitMaxTries = FALSE, 
                       errorHitWallTime = FALSE,
                       userVars = userVars,
                       rules = rules,
                       interventions = interventions)

plot(atex2, show = "genotypes", type = "line")
\end{R}



% Fin del documento
\end{document}
